\section{Conclusion}

This project successfully detailed the design, verification, and empirical synthesis of a highly parameterized 2D Discrete Cosine Transform accelerator. By mapping the computationally dense DCT algorithm into custom SystemVerilog, the project demonstrated the fundamental tradeoffs inherent in digital VLSI design. The hardware-software co-design methodology (utilizing a Python golden reference) allowed for the verification of the mathematical fidelity of the 16-bit Q2.14 fixed-point datapath, consistently achieving an impressive Peak Signal-to-Noise Ratio (PSNR) of 51.05 dB.

Through the systematic evaluation of four distinct architectural variations (D1--D4) on an Intel Cyclone V FPGA, several key lessons were extracted regarding pipelining and spatial parallelism. The baseline D1 architecture confirmed that a serialized MAC engine provides an exceptionally small area footprint (320 ALMs), making it ideal for heavily constrained embedded IoT sensors. However, for real-time high-definition video processing, the D4 configuration proved strictly superior. By orthogonally layering a deep arithmetic pipeline with an 8-way parallel MAC array, the D4 engine achieved a massive throughput of over 3.04 million blocks per second at nearly 200 MHz, easily satisfying the rigorous latency requirements of modern media standards.

\subsection{Future Work}
While this accelerator effectively targets the most computationally expensive stage of image compression, it operates as a standalone IP block. Future work should focus on integrating this module into a complete JPEG or MPEG hardware pipeline. Additionally, migrating the architecture from an FPGA target to an ASIC standard-cell flow would allow for more rigorous power-delay product (PDP) analysis at the standard cell level.
